%% =====================================================================
%%  4e de couverture — Anglais, tome 1 (toujours sur une page paire)
%% =====================================================================
\clearpage\ifodd\value{page}\otPageNotes\clearpage\fi
\thispagestyle{empty}
\begin{tikzpicture}[remember picture,overlay]
\begin{scope}[shift={(current page.south west)}]
  \fill[otNavy] (0,0) rectangle (210mm,297mm);
  \draw[white,opacity=0.05,line width=0.3pt,step=7mm] (0,0) grid (210mm,297mm);
  \fill[otGreen] (0,0) rectangle (210mm,2.4mm);

  \node[anchor=north west,inner sep=0pt,text width=165mm,text=white,font=\otBold\fontsize{22}{27}\selectfont]
    at (24mm,262mm) {Vous partez de zéro ?\\Cette méthode est faite pour vous.};
  \fill[otGreen] (24mm,236mm) rectangle (54mm,237.4mm);
  \node[anchor=north west,inner sep=0pt,text width=150mm,text=white,opacity=0.9,
    font=\otRegular\fontsize{11}{16}\selectfont] at (24mm,228mm)
    {Avec Léa, une Française qui s'installe à New York, apprenez l'anglais dont vous avez
     réellement besoin : vous présenter, parler de votre quotidien, vous débrouiller en ville,
     inviter, raconter. Chaque mot est expliqué en français, chaque son est enregistré,
     chaque exercice est corrigé.};

  % Trois arguments
  \foreach \n/\ico/\titre/\txt in {0/bulle/Une méthode complète/20 unités : dialogues\, vocabulaire\, grammaire\, prononciation\, culture et exercices.,
      1/casque/L'audio partout/Plus de 250 pistes à scanner avec votre téléphone\, en version lente et normale.,
      2/blocs/Pensée pour les francophones/Les pièges du français anticipés\, et des corrigés qui expliquent chaque réponse.}{%
    \pgfmathsetmacro{\x}{24+\n*55}
    \fill[white,opacity=0.07,rounded corners=2mm] (\x mm,140mm) rectangle ++(51mm,48mm);
    \fill[otGreen,rounded corners=1.6mm] (\x mm+4mm,174mm) rectangle ++(10mm,10mm);
    \otDessinIcone{\ico}{white}{otGreen}{6.4mm}{(\x mm+9mm,179mm)}
    \node[anchor=north west,inner sep=0pt,text width=44mm,text=white,align=left] at (\x mm+4mm,170mm)
      {{\otSemi\fontsize{10.4}{12.4}\selectfont\titre}\par\vspace{1.4mm}
       {\otRegular\fontsize{8.8}{11.4}\selectfont\color{white!80!otNavy}\txt}};}

  % Au programme
  \node[anchor=north west,inner sep=0pt,text width=162mm,text=white,align=left] at (24mm,128mm)
    {{\otSemi\fontsize{7.5}{9}\selectfont\color{otGreen}\otCaps[12]{Au programme}}\par\vspace{2.4mm}
     {\otRegular\fontsize{10}{14}\selectfont Saluer et se présenter\enspace·\enspace Épeler, compter, donner son numéro\enspace·\enspace
      Présenter sa famille\enspace·\enspace Raconter sa journée\enspace·\enspace Dire ce qu'on aime\enspace·\enspace
      Demander son chemin\enspace·\enspace Faire des achats\enspace·\enspace Commander au restaurant\enspace·\enspace
      Parler du temps qu'il fait\enspace·\enspace Inviter\enspace·\enspace Raconter un voyage}};

  % Échelle des niveaux
  \foreach \n/\niv in {0/A1,1/A2,2/B1,3/B2,4/C1}{%
    \pgfmathsetmacro{\x}{24+\n*20}
    \ifnum\n=0 \fill[otGreen,rounded corners=1mm] (\x mm,72mm) rectangle ++(18mm,9mm);
    \else \fill[white,opacity=0.12,rounded corners=1mm] (\x mm,72mm) rectangle ++(18mm,9mm);\fi
    \node[text=white,font=\otBold\fontsize{10}{11}\selectfont] at (\x mm+9mm,76.5mm) {\niv};
    \node[anchor=north,text=white,opacity=0.7,font=\otRegular\fontsize{6.8}{8}\selectfont] at (\x mm+9mm,70mm) {Tome \the\numexpr\n+1\relax};}
  \node[anchor=west,text=white,opacity=0.85,font=\otRegular\fontsize{9}{11}\selectfont,text width=60mm] at (128mm,76mm)
    {Une collection en cinq tomes, du niveau débutant au niveau C1.};

  % Pied : logo, site, emplacement du code-barres
  \node[anchor=south west,inner sep=0pt] at (24mm,24mm) {\includegraphics[width=44mm]{logo-blanc}};
  \node[anchor=base west,inner sep=0pt,text=white,opacity=0.75,font=\otMedium\fontsize{9.5}{11}\selectfont]
    at (24mm,16mm) {\otSite};
  \fill[white] (146mm,16mm) rectangle (186mm,40mm);
  \node[text=otMuted,font=\otRegular\fontsize{7}{9}\selectfont,align=center] at (166mm,28mm) {ISBN\\à paraître};
\end{scope}
\end{tikzpicture}
